La realización de este TFG supone la creación de una aplicación web que permite a los usuarios gestionar sus títulos de videojuegos 
favoritos y compartir sus opiniones con otros usuarios. 

Con la ausencia de aplicaciones similares, se ha ido desarrollando este proyecto, estableciendo desde el inicio 
las necesidades y funcionalidades que se querían implementar. Además, se ha realizado un análisis de las 
tecnologías más adecuadas para su desarrollo, teniendo en cuenta la 
experiencia previa en los lenguajes de programación. 

Por otro lado, la realización de mockups ha permitido definir la interfaz de usuario y la experiencia de usuario,
lo que ha facilitado la implementación de la aplicación. La realización de pruebas ha permitido detectar errores y 
mejorar la usabilidad de la aplicación, asegurando que cumple con los requisitos establecidos.

Además, la utilización de Angular y Django ha permitido una implementación más rápida y eficiente. Angular ha proporcionado 
una arquitectura modular y herramientas para la creación de componentes reutilizables(como en el caso de la lista 
de videojuegos). Aunque ha presentado problemas en el desarrollo como el mapeo de los datos de la API a los 
componentes o la comunicación entre componentes, se ha logrado una interfaz de usuario intuitiva y funcional. 
Por otro lado, Django ha facilitado la creación de la API, la gestión de la base de datos y la autenticación de 
usuarios.


En conclusión, la realización de este TFG ha permitido afianzar conocimientos en el desarrollo de aplicaciones web. 
A nivel personal, ha supuesto ha supuesto un apredizaje en la estructuración de proyectos, planificación y 
resolución de problemas y, además, ha resultado espcialmente motivador dar vida 
a una idea pensada desde el principio para los amantes de los videojuegos. 